% !TeX root = ../thuthesis-example.tex

\chapter{总结与展望}

\section{项目总结}

本项目设计并实现了一个基于 Spring Boot 3.2.5 与 Spring Cloud 2023.0.1.0 生态技术栈的水果生鲜运营平台。该系统针对传统单体生鲜电商在部署效率、水平扩展、故障隔离以及多模块维护方面的痛点进行了微服务化改造。系统最终落地为一个由统一 API 网关进行请求分发、以 Nacos 为注册发现与配置管理中心、通过 OpenFeign 进行跨微服务声明式强类型调用、引入 RabbitMQ 消息中间件进行异步解耦、采用本地 MinIO 进行静态图片对象存储、并依托 WebSocket 提供长连接实时推送的高可用分布式系统。

在开发产出方面，本人围绕架构设计、核心业务开发、安全支付隔离、实时消息通知以及系统质量保障完成了完整闭环。需求与设计阶段，本人完成微服务边界定义、网关路由断言、动态跨域配置、Nacos 配置中心架构以及基于 PM2 的长运行进程守护编排。编码阶段，本人对核心业务服务中的分类管理、品种维护及订单持久化接口进行了重构，设计了基于 Spring AOP 的公共字段自动填充切面，并完成支付回调、OpenFeign 客户端、RabbitMQ 队列绑定、WebSocket 长连接广播推送等关键功能。测试阶段，本人利用 Newman 编写了覆盖系统核心功能与网关路由的接口自动化测试集，并通过 Playwright 端到端测试验证主要业务链路。上述工作保证了项目在有限时间内完成高质量交付，也体现了本人对分布式应用开发流程的系统掌握。

\section{遇到的主要技术挑战与解决策略}

在项目的分布式重构与本地多服务联调期间，本人遭遇了数次由于依赖冲突、端口占用、安全拦截以及类型转换引发的底层故障，并通过严谨的代码审计和技术验证逐一予以妥善解决。

第一是关于分布式依赖冲突与 Spring Boot 版本回退的解决。在微服务搭建初期，由于误用了处于里程碑阶段且与若依框架及 MyBatis-Plus 存在严重 API 兼容问题的 Spring Boot 4.x 构建，导致系统编译失败。本人通过将项目基础依赖降级至标准的 Spring Boot 3.2.5 与 Spring Cloud 2023.0.1.0 稳定版，并排除掉冲突的间接依赖包，还原了数据源自动配置类的 Jakarta 命名空间包导入，彻底消除了启动崩溃。此外，针对业务 AOP 字段填充切面在启动时因残留的历史包路径导致切点解析失败的故障，本人通过使用编辑器全局检索，将所有切点配置统一规范为本系统的业务包路径，保证了 Spring 容器的正常启动与切面切点的正确植入。

第二是关于本地端口资源冲突与前端端口自动漂移的解决。在本地开发调试阶段，由于本地腾讯 QQ 进程占用了默认的 8083 端口，导致业务微服务无法正常绑定监听，控制台抛出端口占用异常。本人通过代码重构，将业务微服务迁移至 8088 端口，并同步更新了 AOP 切点拦截范围与各跨服务调用的接口配置。同时，在运行 Playwright 端端测试时，前端预览服务在端口冲突时会自动发生端口漂移，导致自动化测试脚本中硬编码的预览地址失效。本人通过在前端预览配置文件中将预览端口固定为 8087 并开启 strictPort 强制锁定选项，使得 Vite 遇到端口占用时直接报错退出而不是漂移，从而确保了自动化测试路径的绝对一致与可靠。

第三是关于安全过滤链排除与监控端点白名单适配的解决。支付和通知微服务在集成监控模块后，因为继承了若依默认的安全配置，导致外部测试工具发起的 Mock 支付请求、WebSocket 连接以及 Actuator 监控拉取被 Spring Security 拦截并返回 401 未授权错误。本人的解决策略是在两个子服务的启动类上通过 Exclude 属性排除了安全与监控安全自动配置类，完全移除了其安全过滤链拦截。对于必须保留安全拦截的管理服务，本人在 SecurityConfig 配置类中将监控路径添加至匿名白名单，从而打通了本地联调通道与外部监控拉取的连接，保障了测试的顺畅执行。

第四是关于微信异步回调下的 Feign 空上下文及类型转换崩溃的解决。在支付成功异步回调阶段，由于该 HTTP 请求属于微信服务发起，不携带当前登录用户的 JWT 令牌上下文，这导致经过网关和 OpenFeign 转发时因缺少认证头而频繁抛出鉴权失败异常。同时，在订单状态更新接口中，传入的订单号由于带有特定前缀，导致直接调用 Long.valueOf 方法将其转为主键 ID 时触发了 Java 运行时的数字格式化异常。本人通过在业务层引入正则表达式校验，仅对全数字订单号调用转换方法，对含有非数字字符的订单号降级为根据订单字段的多条件查询；同时本人在内部通信的安全策略上进行了局部放行，成功解决了这一跨微服务状态流转过程中的深层异常。

\section{系统局限性与不足}

尽管系统当前作为课程设计项目已经达成了各项技术指标和测试要求，但在应对真正生产级高并发、高可用环境的标准时，依然存在两点主要的局限性。

首先是单物理数据库分表隔离带来的局限。为简化本地运行验证的部署门槛和数据初始化流程，当前各个微服务模块依然共享同一个物理 MySQL 数据库实例，只是通过表前缀进行了逻辑划分。这种共享物理数据源的设计，无法模拟出真实的跨网络物理数据库隔离场景，也使得服务之间的解耦不够彻底。在真正的分布式环境下，每个微服务应拥有完全独立的数据库实例。当前的单库设计避开了分布式事务和分布式锁的复杂场景，对于跨微服务的强一致性数据操作，未能在底层数据库层面完全暴露出分布式一致性问题（例如多物理库提交失败时的数据回滚），这是当前系统最大的局限所在。

其次是缺乏完整的分布式链路追踪与系统级限流手段。为了防止本地开发机器因内存资源不足而导致系统卡顿甚至蓝屏，本人在微服务集成时暂时移除了 Sentinel 限流降级服务以及 SkyWalking 链路追踪服务。这导致系统在面临高并发流量冲击时缺乏流量塑形和网关防刷能力，无法在网关层对恶意流量进行阻断。同时，由于缺乏 SkyWalking 的 APM 监控，微服务之间的多级 Feign 调用链路对运维人员来说是一个黑盒，若某个调用环节发生响应慢、数据库死锁或网络丢包，运维人员无法在可视化大屏上通过调用拓扑图和链路追踪数据对慢 SQL 或故障节点进行快速定位，给后期的系统级调优和维护带来不便。

\section{后续展望}

在未来的系统迭代和技术演进中，本人计划从容器化部署、分布式事务控制以及系统监控治理三个方向对生鲜系统进行深度拓展。

第一是全面实现云原生的容器化部署。后续本人将在 deploy 分支中编写全套 Dockerfile，将管理后台前端应用、五个后端微服务 JAR 包打包为轻量级的 Docker 镜像。同时编写容器化部署配置，将 Nginx 反向代理、Nacos 服务集群、Redis 哨兵高可用实例、RabbitMQ 镜像队列以及 MySQL 主从复制实例整合在统一的虚拟容器网络中。这不仅能够真正免去运行验证机器对本地安装各类中间件的依赖，还能通过一条命令实现一键拉起与弹性扩缩容，真正达到云原生就绪状态。

第二是引入分布式事务协调器以保障数据最终一致性。针对微服务拆分后涉及多物理数据库的写入场景，本人计划在系统中集成阿里开源的 Seata 分布式事务组件。在核心业务服务与支付服务涉及数据变更的方法上，利用 Seata 提供的全局事务注解，将远程 Feign 调用和本地 Mapper 写入包装在同一个全局事务中。通过 AT（自动事务）模式或 TCC（补偿型事务）模式，在网络抖动或下游服务抛出异常时实现全局数据回滚，从而在微服务多库物理隔离后依然能保障订单与支付状态的强一致性。

第三是搭建基于 Prometheus 与 Grafana 的中心化监控体系。针对目前缺乏链路观测的不足，计划在各微服务中集成 Actuator 的 Prometheus 适配包，暴露标准的指标数据采集端点。在容器网络中部署 Prometheus 时序数据库实例，使其每隔十秒自动拉取各微服务实例的 JVM 内存使用、CPU 负载、垃圾回收耗时以及 Active 线程数。最终在 Grafana 中设计搭建多维度的图形可视化大屏，全天候、直观地监控微服务集群的运行状况，提升微服务集群的整体可观测性与运维敏捷度。

第四是引入面向生鲜零售的个性化推荐能力。当前系统已经积累了用户、购物车、订单明细和商品销量数据，后续可在保障隐私与数据最小化原则的前提下，参考推荐系统经典框架与矩阵分解、隐式反馈、深度匹配、多任务排序等方法，对用户复购水果、热销组合和相似商品进行排序建模\cite{ricci,hu2008collaborative,rendle2009bpr,covington,cheng,he2017ncf}。在推荐列表生成后，还可借鉴最大边际相关性思想，在相关性与多样性之间进行平衡，避免首页商品只集中展示单一高频品类\cite{carbonell1998mmr}。
